\subsection{Eventos}

\subsubsection{Regla: verbo y tiempo verbal}

\textbf{Regla:} Los eventos se nombran en \texttt{PascalCase} con un verbo o frase verbal. El presente progresivo indica que la acción está por ocurrir (\texttt{TurnChanging}) y el pasado indica que ya ocurrió (\texttt{TurnChanged}); no se usan prefijos como \texttt{Before} o \texttt{After} (Microsoft, 2023a).

\textbf{Código correcto:}
\begin{lstlisting}[style=csharpstyle]
public class TurnManager
{
    public event EventHandler? TurnChanged;

    protected virtual void OnTurnChanged(EventArgs e)
    {
        TurnChanged?.Invoke(this, e);
    }
}
\end{lstlisting}

\textbf{Código incorrecto:}
\begin{lstlisting}[style=csharpstyle]
public class TurnManager
{
    public event EventHandler? AfterTurnChange;

    protected virtual void RaiseTurnChanged(EventArgs e)
    {
        AfterTurnChange?.Invoke(this, e);
    }
}
\end{lstlisting}

\subsubsection{Regla: \texttt{EventHandler} y \texttt{EventArgs}}

\textbf{Regla:} Se utiliza \texttt{EventHandler<TEventArgs>} para eventos con datos. La clase de datos termina en \texttt{EventArgs}; los parámetros del manejador se llaman \texttt{sender} y \texttt{e}; el método que eleva el evento comienza con \texttt{On} (Microsoft, 2023b).

\textbf{Excepción:} \texttt{sender} y \texttt{e} son la única excepción admitida a la regla general de nombres descriptivos para parámetros (Sección 2.1): siguen la firma estándar del delegado \texttt{EventHandler} de .NET y no se renombran ni se sustituyen por un descarte.

\textbf{Código correcto:}
\begin{lstlisting}[style=csharpstyle]
public event EventHandler<TurnChangedEventArgs>? TurnChanged;

private void HandleTurnChanged(object? sender, TurnChangedEventArgs e)
{
    UpdateCurrentPlayer(e.CurrentPlayer);
}
\end{lstlisting}

\textbf{Código incorrecto:}
\begin{lstlisting}[style=csharpstyle]
public event Action<Player>? TurnChange;

private void HandleTurnChange(Player player)
{
    UpdateCurrentPlayer(player);
}
\end{lstlisting}
